\section{Quadratic acquisition budgets and receiver memory}
\label{sec:quadraticbudget89}
A width bound can conceal unused coherent acquisition.  The natural
policy quantity is the worst-path sum of squared reservations.  We
first retain the separate finite remainder, then identify the sharp
local order under this resource bound.  None of these quantities
measures the dimension of a quantum memory.

\begin{proposition}[Finite policy profile]\label{prop:policyprofile89}
Fix a legal quadratic tube as in Definition~\ref{def:tangenttube86},
and suppose $Q_jH_jQ_j=0$ for every $j$.  Use the constants
$\alpha,r$ and scale restrictions of Lemma~\ref{lem:adaptiveblock88}.
For a reset policy $\pi$, define the hard path bounds
\[
 N_\pi=\sup_{\zeta}\sum_{h\in\zeta}n(h),\qquad
 Q_\pi=\sup_{\zeta}\sum_{h\in\zeta}n(h)^2,
\]
where $\zeta$ ranges over formal terminal histories, and put
\[
 B_\pi=\sup_{\zeta}\sum_{h\in\zeta}
                  (\alpha n(h)+\sqrt{r n(h)})^2.
\]
For every permitted remainder, the final separation of this policy
satisfies
\begin{equation}\label{eq:policyprofile89}
\begin{split}
 D(\pi;E,F)&\le\min\{2,2s\sqrt{B_\pi}\}\\
 &\le\min\{2,2s(\alpha\sqrt{Q_\pi}+\sqrt{rN_\pi})\}.
\end{split}
\end{equation}
The receiver dimension and all receiver instruments are unrestricted.
The complete fresh acquisition remains conditionally independent of
old receiver memory, including all within-block references and memory.
\end{proposition}
\begin{proof}
At history $h$ the fresh adaptive-block lemma gives a deficit at most
$a_h=s^2(\alpha n(h)+\sqrt{rn(h)})^2/2$.
The term $n(h)rs^2$ is the full tube-to-surrogate channel error before
Bures conversion.  Along every terminal history their sum is at most
$s^2B_\pi/2$.  Conditional fidelity accumulation and
$\|\rho-\sigma\|_1\le2\sqrt{1-f(\rho,\sigma)^2}$ prove the first
bound, with the universal trace cap covering large deficits.
On each path the Euclidean triangle inequality gives
\[
 \left[\sum_h(\alpha n(h)+\sqrt{rn(h)})^2\right]^{1/2}
 \le\alpha\left(\sum_hn(h)^2\right)^{1/2}
                      +\sqrt{r\sum_hn(h)}.
\]
Taking suprema proves the second bound.  No expectation of a path cost
or probability of a history has been substituted for a supremum.
The induction uses positive subnormalized operators and hence applies
also at null histories and to trace-class countable direct sums.
\end{proof}

\begin{definition}[A quadratic reservation budget]\label{def:quadraticbudget89}
For integers $N\ge1$ and $N\le Q\le N^2$, let
$D^{\mathrm{reset}}_{N,Q}(E,F)$ be the supremum over reset protocols
with $N_\pi\le N$ and $Q_\pi\le Q$.  The budgets include every
reserved slot on every formal terminal history, whether or not it is
used internally.  There is no additional width constraint.  The
conditional reset and common-control conventions are those of
Definition~\ref{def:reset88}.
\end{definition}

\begin{theorem}[Sharp quadratic budget law]\label{thm:quadraticbudget89}
Fix $E$, a nonzero Hermitian zero-sum tuple $H$ in the one-sided cone,
and $\Lambda\ge0$.  There exist $c,C,s_0>0$, depending only on
$E,H,\Lambda$, such that for all integers $N\ge1$, $N\le Q\le N^2$,
$0<s\le s_0$, and every legal $F$ with
\[
 \|F-E-sH\|_\Sigma
 :=\sum_j\|F_j-E_j-sH_j\|_{\mathrm{op}}\le\Lambda s^2,
\]
one has $c\omega\le D^{\mathrm{reset}}_{N,Q}(E,F)\le C\omega$, where
\begin{equation}\label{eq:quadraticbudget89}
\omega=
\begin{cases}
 \min\{1,\sqrt N s\},&H\in\operatorname{Ran}\mathcal C_E,\\
 \min\{1,\sqrt Q s\},&Q_jH_jQ_j=0\ (\forall j),\quad
                           \Gamma_E(H)\notin\mathcal W_E,\\
 \min\{1,Ns\},&Q_jH_jQ_j\ne0\ (\exists j).
\end{cases}
\end{equation}
The lower tester is independent of the unknown remainder and can
be chosen with a deterministic profile, no inter-block feedback,
and parallel acquisition within its blocks.  Constants are uniform
in $N,Q,F$, not in the fixed tube.  The theorem asserts local orders,
not exact-distance identities or a classification of memory dimensions.
\end{theorem}
\begin{proof}
For coherent tangents, Proposition~\ref{prop:policyprofile89} gives
\[
 D(\pi;E,F)\le\min\{2,2s(\alpha\sqrt Q+\sqrt{rN})\}
             \le\min\{2,2(\alpha+\sqrt r)s\sqrt Q\}.
\]
For the lower put $b=\lfloor Q/N\rfloor$.  Then $1\le b\le N$,
$Nb\le Q$, and $Nb\ge Q/2$.  Take as many blocks of size $b$ as
possible, followed by the positive remainder when necessary.
Their sum is $N$ and, by Corollary~\ref{cor:profilebudget89},
\[
 Q/4\le Nb/2\le V_{\boldsymbol n}\le Nb\le Q.
\]
The lower construction of Theorem~\ref{thm:allocation89} belongs to
the prescribed budget class and gives the coherent lower.  In
particular, the possible quadratic-budget integrality gaps cost only
an absolute factor; the theorem does not claim to attain every
integer $Q$ exactly.

For regular tangents use the unrestricted adaptive horizontal upper
and the independent-product lower.  The all-singleton profile has
square cost $N\le Q$.  The finite upper still contains
$2\beta_h\sqrt Ns+(\Lambda+K_h)Ns^2$, capped at two; putting
$x=\sqrt Ns$ gives the claimed order.  For support openings the same
singleton profile uses the fixed impossible-label event, and the
ordinary adaptive hybrid bound gives the upper.  This lower is a
classical leakage experiment rather than a coherent phase code.
All lower controls are chosen from the fixed tube and public budgets;
the complete finite $m_r\kappa s^2$ errors were paid in
Theorem~\ref{thm:allocation89} before statistical amplification.
\end{proof}

\begin{proposition}[Certified deviations from reset]\label{prop:approxreset89}
Fix a finite or countably branching reset policy with a hard finite
call bound.  At each acquisition history $h$, replace the ideal
preparation channel $R\mapsto R\otimes\sigma_h$ by a channel
$\widetilde{\mathcal P}_h$ with the same input and output systems.
Assume, including auxiliary references,
\[
 \|\widetilde{\mathcal P}_h-(\operatorname{Id}\otimes\sigma_h)
                      \|_\diamond\le\delta_h,\qquad
 \sup_\zeta\sum_{h\in\zeta}\delta_h\le\varepsilon.
\]
All within-block testers, receiver instruments, and reservation rules
are kept fixed and common.  The implemented final separation obeys
\begin{equation}\label{eq:approxreset89}
 \widetilde D(\pi;E,F)
 \le\min\{2,\ D(\pi;E,F)+2\varepsilon\}.
\end{equation}
For tangential directions one may substitute the right side of
\eqref{eq:policyprofile89} for $D(\pi;E,F)$.
For the same pair of ideal and implemented channels the proof gives the
two-sided estimate $|\widetilde D(\pi;E,F)-D(\pi;E,F)|\le2\varepsilon$.
In particular a lower $c\omega$ remains at least $c\omega/2$ if
$\varepsilon\le c\omega/4$; the condition $\varepsilon=o(\omega)$
suffices along a shrinking-signal sequence.
\end{proposition}
\begin{proof}
For either fixed device, let $\varepsilon_h$ bound the remaining
path sum from a history $h$.  Induct backwards on remaining depth
with arbitrary positive incoming operator $R$.  Replacing the
preparation changes the continued output in trace norm by at most
$\delta_h\tr R$, since every continuation is a channel.  For the
ideal preparation followed by the common block and receiver
instrument, the positive child operators have traces summing to
$\tr R$.  The induction bounds the remaining difference by
$(\varepsilon_h-\delta_h)\tr R$.  Their sum is
$\varepsilon_h\tr R$.  At terminals it is zero; countable sums
converge monotonically in trace.  Thus each hypothesis's normalized
output differs by at most $\varepsilon$.  Applying the triangle
inequality to the two hypotheses proves the result.
\end{proof}

The hypothesis is a uniform channel certificate for a specified
common ideal reset preparation.  Marginal separability, finite
policy arithmetic, or an average calibration error does not supply
it.  Fixed nonzero $\varepsilon$ can dominate a shrinking local
signal; no fixed-error asymptotic robustness is inferred.  The same
argument applies to the particular lower construction when its
preparations satisfy the displayed hypothesis, giving the lower
separation minus $2\varepsilon$.

\subsection{A strict tester-level comparison with classical adaptation}
We now give an explicit comparison with the classically adaptive
model of Ohst--Zhang--Nguyen--Pl\'avala--Quintino
\cite[Definition~23 and (64)]{OZNPQ89}.  For two calls their tester
operators have the form
\begin{equation}\label{eq:catester89}
 T_i=\sum_z R_z\otimes S_{i|z},
\end{equation}
where $R$ and each $S_{\cdot|z}$ are single-call testers.  This is
an entire first-call measurement followed by conditional new
preparation, not retention of an unmeasured first-call quantum
register for a later joint readout.

\begin{proposition}[Retained receiver memory]\label{prop:memorycomparison89}
Every two-call tester of the form \eqref{eq:catester89} is a unit-width
reset tester.  The inclusion is strict at the level of tester
operators, already for two-dimensional inputs and two classical
output labels.  This is not an assertion of a strict optimal-score
gap for every ensemble or for every pair of identical measurement
channels.
\end{proposition}
\begin{proof}
Realize $R$ on the first call, retain only its classical outcome
$z$, then prepare and realize $S_{\cdot|z}$ freshly.  This obeys the
unit-width reset condition and proves inclusion.

For strictness, write $I_1,I_2$ for the two qubit inputs and $Y_1,Y_2$
for their classical output registers.  Prepare independent maximally
entangled states on $I_rR_r$, one per call.  Keep the first reference
and label in the receiver, append the second complete pair freshly,
and make a final joint measurement.  One effect of that measurement
is $\Pi_{R_1R_2}\otimes|00\rangle\langle00|_{Y_1Y_2}$, where
$\Pi=|\Phi^+\rangle\langle\Phi^+|$ and
$|\Phi^+\rangle=(|00\rangle+|11\rangle)/\sqrt2$.
For precision, order factors as $I_1,Y_1,I_2,Y_2$ and take the fully
transposed unnormalized Choi representative
$J^\sharp(E)=\sum_yE_y\otimes|y\rangle\langle y|$, obtained from
$(\operatorname{Id}\otimes\mathcal M_E)(|\Omega\rangle\langle\Omega|)$
by full transpose, with $|\Omega\rangle=|00\rangle+|11\rangle$.
Thus $\tr_YJ^\sharp=I$.  The two normalized pairs produce the receiver
operator $(J^\sharp(E)\otimes J^\sharp(F))^{\mathsf T}/4$.
For any receiver effect $M$, the Born probability is
$\tr[M^{\mathsf T}(J^\sharp(E)\otimes J^\sharp(F))]/4$.
This proves the transpose and factor $1/4$ for arbitrary complex effects,
not merely for the real Bell projector.  Grouping the input factors below
means the canonical permutation of the stated factor order.
In the Choi convention with input marginal $I$, direct contraction
of the two normalized maximally entangled preparations gives the
tester effect
\begin{equation}\label{eq:memorywitness89}
 T_0=\tfrac14\Pi^{\mathsf T}_{I_1I_2}
                      \otimes|00\rangle\langle00|_{Y_1Y_2},
 \qquad T_1=\tfrac14 I_{I_1I_2}\otimes I_{Y_1Y_2}-T_0.
\end{equation}
Both are positive and their sum is the deterministic parallel tester
with input marginal $I_{I_1I_2}/4$.  They are block diagonal in the
classical outputs.  The described realization has reset width one:
no old reference enters the second acquisition.
Indeed $0\preceq\Pi^{\mathsf T}\otimes|00\rangle\langle00|\preceq I$,
so the complementary effect is positive.  The deterministic sum is
$\Xi\otimes I_{Y_2}$ with
$\Xi=I_{I_1}\otimes I_{Y_1}\otimes I_{I_2}/4$ and
$\tr_{I_2}\Xi=(I_{I_1}/2)\otimes I_{Y_1}$; the initial marginal has
trace one.  This verifies both causal tester trace constraints.
The cone of positive separable operators across the complete-call cut
$I_1Y_1\mid I_2Y_2$ is closed: its trace-one slice is the compact convex
hull of the compact set of product states, and trace controls the scaling
in its cone.  Hence the limiting sums mentioned below lie in this same
closed cone.  Partial transpose is used only as a necessary condition
for separability.

Partial transpose on the complete first-call factor $I_1Y_1$ leaves
the output projector unchanged.  Since $\Pi$ is real and its partial
transpose has eigenvalues $1/2,1/2,1/2,-1/2$, the partial transpose of
$T_0$ has an eigenvalue $-1/8$.  Every positive sum
\eqref{eq:catester89} has positive partial transpose across
$I_1Y_1|I_2Y_2$, including limits of such sums.  Hence this effect
cannot occur in a tester of that form.  This proves strictness.
\end{proof}

Consequently the constrained-separability condition on complete
first- and second-call tester factors in \cite[(65)--(68)]{OZNPQ89}
cannot be imposed on the whole unit-reset class.  Reset independence
is a factorization at a specified conditional preparation cut; it
does not imply separability of the eventual tester across call
factors.  The witness is a calibration of model inclusions, not a
new memory-discrimination mechanism or a physical experiment.
